% SOURCE AUTHORITY:
% expose_VIII_body.tex lines 2103--2170; scan indices 307--310;
% source folios 296--299.
% Direct scan comparison: physical PDF pages 308--311 in overlapping
% 1100-dpi bands.  The prose and formulas are fully legible at that
% resolution.  The hand-corrected diagram on physical page 310 was also
% inspected directly in a 5000-dpi targeted crop.
% Transparent source dispositions:
% - 6.9(a) begins with j:V->S but consistently defines U=S-Y; V is corrected
%   to U, and j_*G_{mS} is corrected to j_*G_{mU};
% - the source diagram retains a faint obsolete right arrowhead on the Y--S
%   stroke beneath a heavy hand-added left arrowhead.  The correction gives
%   Y->S.  The mark at U is an open-immersion hook, and the source's terminal
%   period is retained.
% Hard stop before section 7, source line 2171 / scan index 310.

\medskip
\noindent\underline{Remarks 6.8.}\quad a) We shall apply 6.6 and 6.7 when
$S$ is a trait and $D$ is the divisor defined by the inclusion of the closed
point $s$ of $S$ into $S$.  In this case $U$ is the spectrum of the fraction
field $K$ of $S$, and $\overline G$ is the usual Neron model of the
multiplicative group over $K$, introduced by Raynaud.  Notice that even in
this special case, the proof of 6.6 and 6.7 uses 6.5 over bases more general
than traits (such as $P$ and its Cartesian powers).  Nevertheless, in this
case it is enough to know the \underline{existence} of the Neron model
$\overline G$ in the special situation under consideration: 6.5 for base
schemes smooth over $S$, and hence 6.6 and 6.7, then follow formally.  For
us, the introduction of $\overline G$ will chiefly serve as a convenient
intermediate step in studying, in the next section, the problem of extending
extensions and biextensions by the multiplicative group.

\smallskip
\noindent b) A Neron--Raynaud model $\overline G$ of
$\underline G_{mU}$, satisfying the fundamental condition (6.4.2) for every
$S'$ smooth over $S$, can be defined under much more general conditions than
those considered above, with $U$ merely a schematically dense open subset
of a normal locally noetherian scheme $S$.  In this case $\overline G$ is an
fppf \underline{sheaf}, not necessarily representable.  Let $\underline D$
denote the sheaf, on the small etale site of $S$, of divisors on $S$ whose
support is contained in $T=S\setminus U$, and, by abuse of language, also
denote by the same letter its canonical extension to a sheaf on
$(\mathrm{Sch})_{/S}$.  Its sections over a relative scheme
$f:S'\to S$ are the sections of $f_{\mathrm{\acute{e}t}}^*(\underline D)$
over $S'$.  The desired model is an extension of groups
\begin{equation}\tag{6.8.1}
 0\longrightarrow\underline G_{mS}\longrightarrow\overline G
 \longrightarrow\underline D\longrightarrow0.
\end{equation}
The construction of a canonical extension (6.8.1), without any normality
hypothesis on $S$, presents no difficulty when one uses the general
description VII 1.1.6; at the same time one obtains (6.4.2) for $S'$ etale
over $S$.  For the same relation to hold for every $S'$ smooth over $S$, it
is equivalent to verify that the formation of the sheaf $\underline D$
commutes with every smooth base change $f:S'\to S$.  For normal $S$, this
follows from EGA $\mathrm{Err}_{\mathrm{IV}}$ 53, 21.4.9.

\medskip
\noindent 6.9.\quad ``\underline{Self-criticism}'' (observations of
P.~Deligne).\footnote{The author of the present Expose having admitted
defeat, the reader is invited, if necessary, to rewrite sections 5 and 6
according to the accompanying directions.}

Sections 5 and 6 are really rather uninspiring.  One periodically has the
impression that the formalism of $i_*,i^*$ is being redone ``by hand'',
without seeing what is geometric and what is ``general nonsense''.

I propose extracting the ``geometric'' statements a), b), c):

\smallskip
\noindent a) (cf. 6.8 b)) Let $j:U\hookrightarrow S$ be a schematically
dense open subset of a normal noetherian scheme $S$.  The etale sheaf
$\underline G_{mS}$ is then identified with a subsheaf of the etale sheaf
$j_*\underline G_{mU}$.  Let $\underline D$ be the quotient etale sheaf in
the exact sequence
\[
 0\longrightarrow\underline G_{mS}\longrightarrow
 j_*\underline G_{mU}\longrightarrow\underline D\longrightarrow0.
\]
\underline{Then} the formation of $\underline D$ commutes with every smooth
base change (EGA $\mathrm{Err}_{\mathrm{IV}}$ 53, 21.4.9); \underline{moreover},
$R^1j_*\underline G_m=0$ (Theorem 90) if the strictly local rings of $S$ are
factorial.

\smallskip
\noindent b) If $U=S\setminus Y$, where $Y$ is a geometrically unibranch
divisor, and if $i$ is the inclusion of $Y$ into $S$, then
\[
 \underline D=i_*\mathbb Z_Y.
\]

\smallskip
\noindent c) If $Y$ is a closed subscheme of a scheme $X$, and $E$ is a set,
then the etale sheaf $i_*\underline E_Y$ is represented by an etale scheme
over $X$, whose formation commutes with every base change.

\medskip
\noindent\underline{Corollary.}  \underline{Let $S$ be a normal noetherian
scheme, let $Y$ be a geometrically unibranch divisor, and let
$U=S\setminus Y$:}
\begin{center}
\begin{tikzcd}[column sep=huge]
 U \arrow[r, hook, "j"'] & S & Y \arrow[l, "i"']
\end{tikzcd}
\qquad .
\end{center}
\underline{There exists a unique smooth group scheme $\mathcal G_m$ over
$S$ such that after every smooth base change one has
$\mathcal G_m|S_{\mathrm{\acute{e}t}}=j_*(\underline G_m)_U$.}

Indeed, the preceding requirements determine $\mathcal G_m$ as a sheaf, for
the etale topology, on the category of schemes smooth over $S$.  By a) and
b), the sequence
\[
 0\longrightarrow\underline G_m\longrightarrow\mathcal G_m
 \xrightarrow{\ q\ }i_*\mathbb Z_Y\longrightarrow0
\]
is exact.  By c), $i_*\mathbb Z_Y$ is representable.  Finally, the morphism
$q$ is representable.

Use of the preceding ``small smooth site'' should make it possible to deduce
the statements of section 6 formally from a).
